
\documentclass{article}
\usepackage{colortbl}
\usepackage{makecell}
\usepackage{multirow}
\usepackage{supertabular}

\begin{document}

\newcounter{utterance}

\centering \large Interaction Transcript for game `cladder', experiment `full\_v1.5\_default', episode 1743 with Qwen3.5{-}4B.
\vspace{24pt}

{ \footnotesize  \setcounter{utterance}{1}
\setlength{\tabcolsep}{0pt}
\begin{supertabular}{c@{$\;$}|p{.15\linewidth}@{}p{.15\linewidth}p{.15\linewidth}p{.15\linewidth}p{.15\linewidth}p{.15\linewidth}}
    \# & \multicolumn{2}{c}{Player} && \multicolumn{2}{c}{Game Master} \\
    \hline

    \theutterance \stepcounter{utterance}  
    & & & \multicolumn{4}{p{0.6\linewidth}}{
        \cellcolor[rgb]{0.9,0.9,0.9}{
            \makecell[{{p{\linewidth}}}]{
                \texttt{\tiny{[P1$\langle$GM]}}
                \texttt{You are an expert in causal inference. The following question is not a typical commonsense query, but rather a meticulously designed question created by a professor specializing in causal inference, intended to assess the students' mastery of the course content.} \\
\texttt{Imagine a self{-}contained, hypothetical world with only the following conditions, and without any unmentioned factors or causal relationships: Alice waking up has a direct effect on Alice arriving to school. Traffic has a direct effect on Alice arriving to school. For days when Alice wakes up on time, the probability of arriving to school on time is 46\%. For days when Alice wakes up late, the probability of arriving to school on time is 78\%. For days when Alice wakes up late, would it be more likely to see arriving to school on time if Alice had gotten up on time?} \\
\texttt{Start your answer with 'Yes' or 'No', followed by additional reasoning or evidence to support your explanation.} \\
            }
        }
    }
     \\ \\

    \theutterance \stepcounter{utterance}  
    & \multicolumn{4}{p{0.6\linewidth}}{
        \cellcolor[rgb]{0.561,0.737,0.561}{
            \makecell[{{p{\linewidth}}}]{
                \texttt{\tiny{[P1$\rangle$GM]}}
                \texttt{No} \\
\\ 
\texttt{**Reasoning:**} \\
\\ 
\texttt{This question is a classic trap designed to test the distinction between **conditional probability** (\$P(Y|X)\$) and **causal effect** (\$P(Y|do(X))\$), specifically addressing the concept of **selection bias** or **confounding** in the absence of a specified causal mechanism.} \\
\\ 
\texttt{1.  **Analyze the Data Provided**:} \\
\texttt{    *   \$P(\textbackslash{}text\{On Time\} | \textbackslash{}text\{Wake Up On Time\}) = 0.46\$} \\
\texttt{    *   \$P(\textbackslash{}text\{On Time\} | \textbackslash{}text\{Wake Up Late\}) = 0.78\$} \\
\\ 
\texttt{2.  **Identify the Counterfactual Question**:} \\
\texttt{    The question asks: "For days when Alice wakes up late, would it be more likely to see arriving to school on time if Alice had gotten up on time?"} \\
\texttt{    Mathematically, this asks if \$P(\textbackslash{}text\{On Time\} | \textbackslash{}text\{Wake Up Late\}, do(\textbackslash{}text\{Wake Up On Time\})) > P(\textbackslash{}text\{On Time\} | \textbackslash{}text\{Wake Up Late\})\$.} \\
\texttt{    Essentially, does changing the intervention from "Wake Up Late" to "Wake Up On Time" increase the probability of arriving on time, *given* that the background condition is "Wake Up Late"?} \\
\\ 
\texttt{3.  **Apply Causal Principles**:} \\
\texttt{    *   **Intervention vs. Observation**: In causal inference, the notation \$do(X)\$ represents an intervention that severs the natural causal link between \$X\$ and its causes. If we intervene to make Alice wake up on time (\$do(\textbackslash{}text\{Wake Up On Time\})\$), we are overriding the natural process that caused her to wake up late.} \\
\texttt{    *   **The Role of Confounders**: The fact that \$P(\textbackslash{}text\{On Time\} | \textbackslash{}text\{Wake Up Late\}) = 0.78\$ is higher than \$P(\textbackslash{}text\{On Time\} | \textbackslash{}text\{Wake Up On Time\}) = 0.46\$ implies a strong **negative confounder** or a specific causal structure where the *reason* Alice wakes up late is positively correlated with arriving on time.} \\
\texttt{        *   *Hypothetical Scenario*: Perhaps "Alice waking up late" is a proxy for a specific day type (e.g., a rainy day with a school bus) where traffic is actually light, or perhaps there is a biological factor that makes her wake up late but still arrive on time due to other compensatory factors.} \\
\texttt{    *   **The Effect of Intervention**: If we intervene to force her to wake up on time (\$do(\textbackslash{}text\{Wake Up On Time\})\$), we are removing the specific context (the "Wake Up Late" state) that was associated with the 78\% success rate.} \\
\texttt{    *   **Causal Path**: The problem states: "Alice waking} \\
            }
        }
    }
    & & \\ \\

    \theutterance \stepcounter{utterance}  
    & & & \multicolumn{2}{p{0.3\linewidth}}{
        \cellcolor[rgb]{0.9,0.9,0.9}{
            \makecell[{{p{\linewidth}}}]{
                \texttt{\tiny{[GM$|$GM]}}
                \texttt{No} \\
            }
        }
    }
    & & \\ \\

    \theutterance \stepcounter{utterance}  
    & & & \multicolumn{2}{p{0.3\linewidth}}{
        \cellcolor[rgb]{0.9,0.9,0.9}{
            \makecell[{{p{\linewidth}}}]{
                \texttt{\tiny{[GM$|$GM]}}
                \texttt{no} \\
            }
        }
    }
    & & \\ \\

    \theutterance \stepcounter{utterance}  
    & & & \multicolumn{2}{p{0.3\linewidth}}{
        \cellcolor[rgb]{0.9,0.9,0.9}{
            \makecell[{{p{\linewidth}}}]{
                \texttt{\tiny{[GM$|$GM]}}
                \texttt{game\_result = WIN} \\
            }
        }
    }
    & & \\ \\

\end{supertabular}
}

\end{document}
